\section{Introduction}
A multi-task leaderboard is a matrix of scores, one row per model and one column per sub-dataset, plus an aggregation rule that turns each row into a rank.
Rank affects adoption, publication and funding, so developers have a reason to act on it.
One available action is to add data related to some sub-datasets to the training of a particular model, which raises its scores on those sub-datasets \citep{zhou2023dont,yang2023rephrased}.
How much of this does a developer need to put a chosen model first, and does the answer depend on whether the leaderboard averages, takes a median, or counts wins?

We study this as a minimum-cost problem on public score matrices.
Our setting is deliberately simple. One target model may raise its score on any set of tasks, up to a cap. Every other model is fixed. The target is first if it is strictly ahead of every rival under the rule.
The cost is measured in leaderboard points, three ways: tasks touched ($K$), total points added ($L_1$) and largest single-task gain ($L_\infty$), each minimised separately and exactly. The cap is the main modelling assumption, so it is a swept parameter, from the theoretical ceiling to fixed per-task gains and empirical gains of fine-tuned models over their base models.

\paragraph{Contributions.}
(i) A formal model of targeted investment under an aggregation rule, with three cost measures, a cap parameter, and exact solvers for six score rules and three rank rules (Section~\ref{sec:method}). The closed forms for the mean, minimum and median are elementary; we state them to validate the solvers and to say what the rules cost in the simplest case (Appendix~\ref{app:theory}). What is not derivable in a few lines is the exact treatment of the trimmed mean and the rank rules, by mixed-integer programmes, and the results below.
(ii) An evaluation on every non-leading model of two public leaderboards: HELM MMLU \citep{liang2022holistic,hendrycks2021measuring} and the Open LLM Leaderboard, with uncertainty over targets, over evaluation noise, over task subsets and over rule-specific choices (Section~\ref{sec:experiments}). The empirical contribution is the map of how cost and feasibility depend on rule, cost measure and cap, and how a fixed investment fares under evaluation noise.
(iii) Six claims fixed in the plan before the grids, reported with criterion, measured value and verdict, including where a criterion was loose or a claim reworded after the results (Table~\ref{tab:claims}).

\paragraph{Findings.}
\emph{Cost depends on the rule, but not in a single order.} In total points the mean rule is several times dearer than the minimum rule on both leaderboards and the median rule on HELM, but not in the largest single gain, and the number of tasks it needs is among the smallest (Section~\ref{sec:headline}).
\emph{Caps decide whether first place is reachable at all.} With 10 points per task at most \num{ollb.gain10.min.feas}\% of OLLB targets can reach first place under any score rule; the cap at which half of the targets can win is \num{gs.helm.crossmin}--\num{gs.helm.crossmax} points on HELM and \num{gs.ollb.crossmin}--\num{gs.ollb.crossmax} on OLLB, in a different rule order on each (Section~\ref{sec:caps}).
\emph{Task-count intuitions do not generalise}: $K(\text{mean}) \le K(\text{trimmed}) \le K(\text{median}) \le K(\text{min})$ fails for many targets (Section~\ref{sec:order}).
\emph{Rank rules are cheaper in points than the mean when unconstrained; under caps plurality and the Condorcet requirement are feasible for fewer targets and Borda for more} (Section~\ref{sec:ordinal}).
\emph{Near the top the nominal minimum is not enough under noise} (Section~\ref{sec:noise}).

\paragraph{What we do not claim.}
Investment is measured in score points on the leaderboard, not in training data, compute or money: we do not model how many points a given amount of data buys. We study one deviating developer against fixed rivals, on two leaderboards which are both narrow (Section~\ref{sec:discussion}).
